\subsection{Nonzero steps and finite backtracking}

For every formed trial, write $r_k=\norm{s_k}$. The index $k$
continues to count all trials, including rejected ones.

\begin{lemma}[Nonzero trial steps]\label{lem:nonzero}
Every trial formed by Algorithm~\ref{alg:utr} has $h_k>0$,
$R_k>0$, and $r_k>0$.
\end{lemma}
\begin{proof}
If $h_k=0$, both terms in its defining maximum vanish. This gives
$\nabla f(x_k)=0$ and a positive semidefinite Hessian, so the stopping
test would have succeeded before the trial was formed.
Thus $h_k>0$, and $R_k>0$ follows from $\sigma_k>0$.

If $\nabla f(x_k)\ne0$, the first equation in \eqref{eq:kkt}
rules out $s_k=0$. If $\nabla f(x_k)=0$, failure of the stopping
test implies $\lambda_{\min}(\nabla^2f(x_k))<-\delta<0$.
The scale definition then gives
\[
 2\sqrt{\sigma_kh_k}
 =-\tfrac14\lambda_{\min}(\nabla^2f(x_k)).
\]
Consequently the smallest eigenvalue of the model Hessian is
$\tfrac34\lambda_{\min}(\nabla^2f(x_k))<0$. A sufficiently
small nonzero vector in a corresponding eigenvector direction is
feasible and has strictly negative model value, whereas the model
value at zero is zero. Hence zero is not a global minimizer.
In either case $r_k>0$.
\end{proof}

\begin{lemma}[Safe trial scale]\label{lem:safe}
Both acceptance tests hold for trial $k$ whenever
\begin{equation}\label{eq:safe}
 \sigma_k\ge L_0+L_1\norm{\nabla f(x_k)}
        +\sqrt{L_1}[-\lambda_{\min}(\nabla^2 f(x_k))]_+.
\end{equation}
Thus the number of consecutive rejected trials at any fixed point
is finite.
\end{lemma}
\begin{proof}
Taking the inner product of the first KKT equation with $s_k$
and using Cauchy--Schwarz gives
\[
 \begin{aligned}
 (2\sqrt{\sigma_kh_k}+\lambda_k)r_k^2
 &=-\nabla f(x_k)^Ts_k-s_k^T\nabla^2f(x_k)s_k\\
 &\le\norm{\nabla f(x_k)}r_k
       +[-\lambda_{\min}(\nabla^2f(x_k))]_+r_k^2.
 \end{aligned}
\]
By Lemma~\ref{lem:nonzero}, division by $r_k$ is legitimate.
The same KKT equation identifies the linearized gradient, so
\begin{equation}\label{eq:linear-gradient}
 \begin{aligned}
 (2\sqrt{\sigma_kh_k}+\lambda_k)r_k
 &\le\norm{\nabla f(x_k)}
       +[-\lambda_{\min}(\nabla^2f(x_k))]_+r_k,\\
 \norm{\nabla f(x_k)+\nabla^2f(x_k)s_k}
 &=(2\sqrt{\sigma_kh_k}+\lambda_k)r_k.
 \end{aligned}
\end{equation}
For $u\in[0,1]$, the identity
\[
 \nabla f(x_k)+u\nabla^2f(x_k)s_k
 =(1-u)\nabla f(x_k)
    +u\bigl(\nabla f(x_k)+\nabla^2f(x_k)s_k\bigr)
\]
and the triangle inequality therefore imply
\begin{equation}\label{eq:linear-segment}
 \norm{\nabla f(x_k)+u\nabla^2f(x_k)s_k}
 \le\norm{\nabla f(x_k)}
       +[-\lambda_{\min}(\nabla^2f(x_k))]_+r_k.
\end{equation}

Under \eqref{eq:safe}, each nonnegative term on its right-hand
side is at most $\sigma_k$. In particular,
\[
 \begin{aligned}
 L_1h_k
 &=\max\left\{L_1\norm{\nabla f(x_k)},
       \frac{L_1[-\lambda_{\min}(\nabla^2f(x_k))]_+^2}
            {64\sigma_k}\right\}\\
 &\le\max\{\sigma_k,\sigma_k/64\}=\sigma_k.
 \end{aligned}
\]
Since $r_k\le R_k$, it follows that
\begin{equation}\label{eq:safe-radius}
 \sqrt{L_1}r_k\le\sqrt{L_1}R_k
   =\frac14\sqrt{\frac{L_1h_k}{\sigma_k}}\le\frac14,
 \qquad L_1r_k^2\le\frac1{16}.
\end{equation}
These inequalities also hold when $L_1=0$ and do not divide by
$L_1$.

Continuity of $\nabla^3f$ on the compact trial segment ensures
that
\[
 M_k:=\max_{0\le u\le1}\norm{\nabla^3f(x_k+us_k)}<\infty.
\]
Taylor's integral formula for the gradient is
\[
 \begin{aligned}
 \nabla f(x_k+us_k)
 &=\nabla f(x_k)+u\nabla^2f(x_k)s_k\\
 &\quad+\int_0^u(u-\theta)
          \nabla^3f(x_k+\theta s_k)[s_k,s_k]\,d\theta.
 \end{aligned}
\]
Here the third derivative applied to two vectors is viewed as
a vector through the Euclidean inner product. Its integral has
norm at most
$M_kr_k^2\int_0^u(u-\theta)\,d\theta
 =\tfrac12u^2M_kr_k^2$. Combining this with
\eqref{eq:linear-segment} yields, for every $u\in[0,1]$,
\[
 \norm{\nabla f(x_k+us_k)}
 \le\norm{\nabla f(x_k)}
       +[-\lambda_{\min}(\nabla^2f(x_k))]_+r_k
       +\tfrac12M_kr_k^2.
\]
Apply the generalized smoothness assumption at every point of
the segment and then take the maximum. Using
\eqref{eq:safe-radius} and \eqref{eq:safe} gives
\[
 \begin{aligned}
 M_k
 &\le L_0+L_1\norm{\nabla f(x_k)}
       +L_1[-\lambda_{\min}(\nabla^2f(x_k))]_+r_k
       +\tfrac12L_1M_kr_k^2\\
 &\le L_0+L_1\norm{\nabla f(x_k)}
       +\tfrac14\sqrt{L_1}[-\lambda_{\min}(\nabla^2f(x_k))]_+
       +\tfrac1{32}M_k\\
 &\le\sigma_k+\tfrac1{32}M_k.
 \end{aligned}
\]
Subtracting $M_k/32$ proves $M_k\le32\sigma_k/31$.

The scalar Taylor formula similarly gives
\[
 \begin{aligned}
 &f(x_k+s_k)-f(x_k)-\nabla f(x_k)^Ts_k
                    -\tfrac12s_k^T\nabla^2f(x_k)s_k\\
 &\qquad=\frac12\int_0^1(1-\theta)^2
     \nabla^3f(x_k+\theta s_k)[s_k,s_k,s_k]\,d\theta.
 \end{aligned}
\]
Its absolute value is at most
$M_kr_k^3/6\le(16/93)\sigma_kr_k^3\le\sigma_kr_k^3$.
The gradient remainder at $u=1$ has norm at most
$M_kr_k^2/2\le(16/31)\sigma_kr_k^2\le3\sigma_kr_k^2$.

The first KKT equation and the positive semidefinite condition
in \eqref{eq:kkt} imply
\[
 \begin{aligned}
 &-\nabla f(x_k)^Ts_k-\tfrac12s_k^T\nabla^2f(x_k)s_k\\
 &=\tfrac12s_k^T\bigl(\nabla^2f(x_k)
              +(2\sqrt{\sigma_kh_k}+\lambda_k)I\bigr)s_k
       +\tfrac12(2\sqrt{\sigma_kh_k}+\lambda_k)r_k^2\\
 &\ge\tfrac12(2\sqrt{\sigma_kh_k}+\lambda_k)r_k^2.
 \end{aligned}
\]
Since $\sigma_kr_k\le\sigma_kR_k=\sqrt{\sigma_kh_k}/4$,
subtracting the scalar Taylor remainder proves
\[
 \begin{aligned}
 f(x_k)-f(x_k+s_k)
 &\ge\left(\sqrt{\sigma_kh_k}+\tfrac12\lambda_k\right)r_k^2
          -\sigma_kr_k^3\\
 &\ge\tfrac34\sqrt{\sigma_kh_k}\,r_k^2
          +\tfrac12\lambda_kr_k^2.
 \end{aligned}
\]
This is \eqref{eq:accept-value}. For the gradient test, the
remainder bound and \eqref{eq:linear-gradient} give
\[
 \begin{aligned}
 \norm{\nabla f(x_k+s_k)}
 &\le(2\sqrt{\sigma_kh_k}+\lambda_k)r_k+3\sigma_kr_k^2\\
 &\le2\sqrt{\sigma_kh_k}R_k+3\sigma_kR_k^2+\lambda_kr_k\\
 &=\tfrac12h_k+\tfrac3{16}h_k+\lambda_kr_k
   =\tfrac{11}{16}h_k+\lambda_kr_k.
 \end{aligned}
\]
Thus \eqref{eq:accept-gradient} holds as well. During consecutive
rejections the point $x_k$ stays fixed, so the right-hand side
of \eqref{eq:safe} is a fixed finite number. Each rejection
doubles the positive trial scale; finitely many doublings reach
that number. The resulting trial must be accepted.
\end{proof}

\subsection{Accepted trials, descent, and step geometry}

Consider any finite prefix ending immediately after its $K$th
accepted trial, where $K\ge1$. Let
$\tau_0<\cdots<\tau_{K-1}$ be its accepted trial indices, and
define the endpoint index $\tau_K=\tau_{K-1}+1$. This last index
denotes the point after the final acceptance; it need not be an
accepted trial. Rejections do not change the iterate, and therefore
\begin{equation}\label{eq:accepted-points}
 x_{\tau_0}=x_0,\qquad
 x_{\tau_{j+1}}=x_{\tau_j}+s_{\tau_j}
 \quad(0\le j<K).
\end{equation}
If the prefix is extended, the endpoint index can change to the
next accepted index, but its associated point remains the same
through all intervening rejections.

For the scale recurrence only, use the formal convention
$\sigma_{\tau_{-1}}:=2\sigma_0$; this does not introduce an
additional trial. Then, for $0\le j<K$,
\begin{equation}\label{eq:backtrack-upper}
 \begin{aligned}
 \sigma_{\tau_j}\le\max\bigl\{\sigma_{\tau_{j-1}}/2,\,
 &2\bigl(L_0+L_1\norm{\nabla f(x_{\tau_j})}\\
 &\hspace{1em}+\sqrt{L_1}
       [-\lambda_{\min}(\nabla^2f(x_{\tau_j}))]_+\bigr)\bigr\}.
 \end{aligned}
\end{equation}
Indeed, the first trial at this point uses
$\sigma_{\tau_{j-1}}/2$, including $j=0$ under the convention.
If this trial succeeds, the first term in the maximum suffices.
Otherwise the trial immediately preceding $\tau_j$ was rejected,
at the same point and at scale $\sigma_{\tau_j}/2$.
By Lemma~\ref{lem:safe}, this rejected scale is strictly smaller
than the safe threshold in \eqref{eq:safe}. Multiplying that
strict inequality by two proves the second alternative in
\eqref{eq:backtrack-upper}.

For every accepted trial $k$, its acceptance test, $r_k>0$, and
$\sqrt{\sigma_kh_k}\ge4\sigma_kr_k$ show
\[
 f(x_k)-f(x_{k+1})
 \ge\tfrac34\sqrt{\sigma_kh_k}\,r_k^2
       +\tfrac12\lambda_kr_k^2
 \ge3\sigma_kr_k^3>0.
\]
Using \eqref{eq:accepted-points}, these inequalities telescope:
\begin{equation}\label{eq:descent-budget}
 \begin{gathered}
 f(x_{\tau_j})-f(x_{\tau_{j+1}})
 \ge3\sigma_{\tau_j}r_{\tau_j}^3>0
 \quad(0\le j<K),\\
 \begin{aligned}
 3\sum_{j=0}^{K-1}\sigma_{\tau_j}r_{\tau_j}^3
 &\le\sum_{j=0}^{K-1}
       \bigl(f(x_{\tau_j})-f(x_{\tau_{j+1}})\bigr)\\
 &=f(x_0)-f(x_{\tau_K})\le f(x_0)-f_* =\Delta.
 \end{aligned}
 \end{gathered}
\end{equation}
Each summand is nonnegative, so the same upper bound applies to
any subinterval of accepted trials. No limit of iterates or
infinite sum is needed for this argument.

\begin{lemma}[Boundary and interior steps]\label{lem:geometry}
Put $\rho=11/16$. An accepted boundary trial $k$, for which
$r_k=R_k$, satisfies
\begin{equation}\label{eq:boundary}
 \begin{aligned}
 \norm{\nabla f(x_k)}&\le16\sigma_kr_k^2,\\
 \norm{\nabla f(x_{k+1})}&\le51\sigma_kr_k^2,\\
 [-\lambda_{\min}(\nabla^2f(x_k))]_+&\le32\sigma_kr_k,\qquad
 \lambda_k\le10\sqrt{\sigma_kh_k}.
 \end{aligned}
\end{equation}
An accepted interior trial $k$, for which $r_k<R_k$, satisfies
\begin{equation}\label{eq:interior}
 \begin{aligned}
 h_k&=\norm{\nabla f(x_k)},\qquad\lambda_k=0,\\
 \norm{\nabla f(x_{k+1})}&\le\rho\norm{\nabla f(x_k)},\\
 [-\lambda_{\min}(\nabla^2f(x_k))]_+
 &\le2\sqrt{\sigma_k\norm{\nabla f(x_k)}}.
 \end{aligned}
\end{equation}
\end{lemma}
\begin{proof}
The definition of $h_k$ implies both
$\norm{\nabla f(x_k)}\le h_k$ and
$[-\lambda_{\min}(\nabla^2f(x_k))]_+
 \le8\sqrt{\sigma_kh_k}$, without requiring a safe scale.
If $r_k=R_k$, then
\[
 h_k=16\sigma_kr_k^2,\qquad
 \sqrt{\sigma_kh_k}=4\sigma_kr_k.
\]
These identities give the first and third bounds in
\eqref{eq:boundary}. The KKT inequality
\eqref{eq:linear-gradient}, which was derived before imposing
\eqref{eq:safe}, gives
\[
 \begin{aligned}
 \lambda_k
 &\le\frac{\norm{\nabla f(x_k)}}{R_k}
        +[-\lambda_{\min}(\nabla^2f(x_k))]_+
        -2\sqrt{\sigma_kh_k}\\
 &\le4\sqrt{\sigma_kh_k}+8\sqrt{\sigma_kh_k}
        -2\sqrt{\sigma_kh_k}
   =10\sqrt{\sigma_kh_k}.
 \end{aligned}
\]
The gradient acceptance test now yields
\[
 \begin{aligned}
 \norm{\nabla f(x_{k+1})}
 &\le\tfrac{11}{16}h_k+\lambda_kr_k\\
 &\le\tfrac{11}{16}(16\sigma_kr_k^2)
       +10(4\sigma_kr_k)r_k
   =51\sigma_kr_k^2.
 \end{aligned}
\]

If $r_k<R_k$, complementarity implies $\lambda_k=0$.
The positive semidefinite condition then gives
\[
 [-\lambda_{\min}(\nabla^2f(x_k))]_+
 \le2\sqrt{\sigma_kh_k}.
\]
If the curvature term were strictly larger than the gradient
term in the definition of $h_k$, it would be positive and give
\[
 2\sqrt{\sigma_kh_k}
 =\tfrac14[-\lambda_{\min}(\nabla^2f(x_k))]_+>0.
\]
The two preceding displays would force a positive number to be
at most one quarter of itself, which is impossible. Therefore
$h_k=\norm{\nabla f(x_k)}$. Substitution into the semidefinite
bound proves the last inequality of \eqref{eq:interior};
substitution into the gradient acceptance test, with
$\lambda_k=0$, proves its contraction inequality.
\end{proof}
